\documentclass[12pt, a4paper]{article}

\usepackage[utf8]{inputenc}
\usepackage[T1]{fontenc}
\usepackage{geometry}
\usepackage{xcolor}
\usepackage{listings}
\usepackage{hyperref}
\usepackage{fancyhdr}
\usepackage{parskip}
\usepackage{microtype}
\usepackage{titlesec}
\usepackage{amsmath}
\usepackage{booktabs}
\usepackage{array}

\geometry{a4paper, top=2.5cm, bottom=2.5cm, left=3cm, right=2.5cm}
\setlength{\headheight}{15pt}

% ─────────────────────────────────────────────────────────────────────────────
%  Paleta VSCode Dark+
% ─────────────────────────────────────────────────────────────────────────────
\definecolor{vscBg}      {HTML}{1E1E1E}
\definecolor{vscFg}      {HTML}{D4D4D4}
\definecolor{vscKeyword} {HTML}{569CD6}
\definecolor{vscType}    {HTML}{4EC9B0}
\definecolor{vscString}  {HTML}{CE9178}
\definecolor{vscComment} {HTML}{6A9955}
\definecolor{vscNumber}  {HTML}{B5CEA8}
\definecolor{vscLineNum} {HTML}{858585}
\definecolor{vscAccent}  {HTML}{DCDCAA}
\definecolor{vscBorder}  {HTML}{404040}

% ─────────────────────────────────────────────────────────────────────────────
%  Estilo listings: VSCode Dark+
% ─────────────────────────────────────────────────────────────────────────────
\lstdefinestyle{vscode}{
  backgroundcolor  = \color{vscBg},
  basicstyle       = \ttfamily\footnotesize\color{vscFg},
  keywordstyle     = \color{vscKeyword}\bfseries,
  stringstyle      = \color{vscString},
  commentstyle     = \color{vscComment}\itshape,
  numberstyle      = \ttfamily\scriptsize\color{vscLineNum},
  identifierstyle  = \color{vscFg},
  numbers          = left,
  numbersep        = 12pt,
  tabsize          = 4,
  showspaces       = false,
  showstringspaces = false,
  breaklines       = true,
  breakatwhitespace= false,
  frame            = single,
  frameround       = tttt,
  framerule        = 0.5pt,
  rulecolor        = \color{vscBorder},
  xleftmargin      = 2.2em,
  framexleftmargin = 1.8em,
  keepspaces       = true,
  captionpos       = b,
  aboveskip        = 1em,
  belowskip        = 0.5em,
  literate=
    {á}{{\'{a}}}1 {é}{{\'{e}}}1 {í}{{\'{i}}}1 {ó}{{\'{o}}}1 {ú}{{\'{u}}}1
    {ñ}{{\~{n}}}1 {Á}{{\'{A}}}1 {É}{{\'{E}}}1 {Í}{{\'{I}}}1 {Ó}{{\'{O}}}1
    {Ú}{{\'{U}}}1 {Ñ}{{\~{N}}}1 {ü}{{\"u}}1   {Ü}{{\"U}}1
    {à}{{\`{a}}}1 {è}{{\`{e}}}1 {â}{{\^{a}}}1 {ê}{{\^{e}}}1
    {î}{{\^{i}}}1 {ô}{{\^{o}}}1 {û}{{\^{u}}}1 {ç}{{\c{c}}}1 {Ç}{{\c{C}}}1,
}

\lstdefinestyle{vsxml}{
  style            = vscode,
  language         = XML,
  keywordstyle     = \color{vscType}\bfseries,
  morekeywords     = {uses-permission, uses-feature, application, activity,
                      service, intent-filter, action, category, property,
                      android, manifest},
}

\lstset{style=vscode, language=Java}

% ─────────────────────────────────────────────────────────────────────────────
%  Encabezado / pie
% ─────────────────────────────────────────────────────────────────────────────
\pagestyle{fancy}
\fancyhf{}
\fancyhead[L]{\small\textsc{Android Helios}}
\fancyhead[R]{\small\textit{\leftmark}}
\fancyfoot[C]{\small\thepage}
\renewcommand{\headrulewidth}{0.4pt}

% ─────────────────────────────────────────────────────────────────────────────
%  Formato de secciones
% ─────────────────────────────────────────────────────────────────────────────
\titleformat{\section}[block]
  {\large\bfseries}{\thesection.}{0.8em}{}
  [\vspace{-0.4em}\rule{\textwidth}{0.4pt}\vspace{0.2em}]

\titleformat{\subsection}[block]
  {\normalsize\bfseries}{\thesubsection.}{0.6em}{}

\titleformat{\subsubsection}[block]
  {\normalsize\itshape}{\thesubsubsection.}{0.6em}{}

\hypersetup{
  colorlinks = true,
  linkcolor  = vscKeyword,
  urlcolor   = vscType,
  citecolor  = vscKeyword,
  pdftitle   = {Android Helios -- Informe Tecnico},
  pdfauthor  = {Facundo Arito},
}

% ═════════════════════════════════════════════════════════════════════════════
\begin{document}
% ═════════════════════════════════════════════════════════════════════════════

% ─── Portada ──────────────────────────────────────────────────────────────────
\begin{titlepage}
  \centering
  \vspace*{3cm}
  {\Huge\bfseries Android Helios\par}
  \vspace{0.6cm}
  {\Large Informe T\'ecnico del Proyecto\par}
  \vspace{2.5cm}
  {\color{vscKeyword}\rule{0.55\textwidth}{1.5pt}}\\[0.8em]
  {\large App Android: Linterna controlada por sacudida\par}
  {\small Aceler\'ometro + CameraManager + Gr\'afico en tiempo real + Foreground Service\par}
  {\color{vscKeyword}\rule{0.55\textwidth}{1.5pt}}
  \vspace{2.5cm}

  \begin{tabular}{@{}ll@{}}
    \textbf{Autor:}       & Facundo Arito \\[0.4em]
    \textbf{Materia:}     & Computaci\'on Aplicada a la Ingenier\'ia \\[0.4em]
    \textbf{Repositorio:} & \texttt{github.com/wopassord/Android\_Helios} \\[0.4em]
    \textbf{Fecha:}       & Junio de 2026 \\
  \end{tabular}

  \vfill
  {\small Stack: Java \textbullet{} Android Studio \textbullet{} Gradle (Kotlin DSL)
  \textbullet{} minSdk 24 \textbullet{} targetSdk 36}
\end{titlepage}

% ─── Índice ───────────────────────────────────────────────────────────────────
\tableofcontents
\clearpage

% ═════════════════════════════════════════════════════════════════════════════
\section{Introducci\'on}
% ═════════════════════════════════════════════════════════════════════════════

Android Helios es una aplicaci\'on Android desarrollada en Java cuyo objetivo
principal es \textbf{encender o apagar la linterna del tel\'efono mediante un
patr\'on de sacudida configurable}. La app integra dos m\'odulos de hardware:

\begin{itemize}
  \item \textbf{Aceler\'ometro:} detecta el movimiento del dispositivo,
        calcula la magnitud de la aceleraci\'on en unidades de gravedad y
        aplica un algoritmo de detecci\'on de sacudidas.
  \item \textbf{C\'amara / Flash:} controla la linterna a trav\'es de
        \texttt{CameraManager.setTorchMode()}, sin necesidad de abrir un
        stream de c\'amara.
\end{itemize}

La interfaz expone tres controles al usuario: un indicador de estado de la
linterna (c\'irculo cromado + texto), un gr\'afico en tiempo real de la
aceleraci\'on medida, y un deslizador de sensibilidad. Un toggle en la esquina
superior derecha permite que la detecci\'on contin\'ue activa incluso cuando el
usuario sale de la aplicaci\'on.

% ═════════════════════════════════════════════════════════════════════════════
\section{Estructura del Proyecto}
% ═════════════════════════════════════════════════════════════════════════════

El proyecto sigue la estructura est\'andar de Android Studio con Gradle (Kotlin
DSL). Los archivos programados son:

\vspace{0.5em}
\begin{center}
\begin{tabular}{@{} l p{9.5cm} @{}}
  \toprule
  \textbf{Archivo} & \textbf{Rol} \\
  \midrule
  \texttt{MainActivity.java}
    & Activity \'unica. Coordina el aceler\'ometro, la linterna, el gr\'afico
      en tiempo real y toda la l\'ogica de UI. \\[0.4em]
  \texttt{AccelerometerGraphView.java}
    & Vista personalizada (\texttt{extends View}) que dibuja el gr\'afico de
      g-force con Canvas de Android. \\[0.4em]
  \texttt{TorchService.java}
    & Foreground Service. Replica la detecci\'on de sacudidas y el toggle de
      la linterna cuando la app est\'a en segundo plano. \\[0.4em]
  \texttt{activity\_main.xml}
    & Layout principal: indicador circular, gr\'afico, slider de sensibilidad
      y toggle de segundo plano. \\[0.4em]
  \texttt{AndroidManifest.xml}
    & Declara permisos (CAMERA, FOREGROUND\_SERVICE) y registra el servicio. \\[0.4em]
  \texttt{drawable/circle\_indicator.xml}
    & Shape drawable oval que sirve de indicador visual del estado de la
      linterna. \\
  \bottomrule
\end{tabular}
\end{center}

% ═════════════════════════════════════════════════════════════════════════════
\section{La Linterna: \texttt{CameraManager}}
% ═════════════════════════════════════════════════════════════════════════════

\subsection{API utilizada}

Android expone el control del flash a trav\'es de
\texttt{android.hardware.camera2.CameraManager}. El m\'etodo
\texttt{setTorchMode(String cameraId, boolean enabled)}, disponible desde
API~23, enciende o apaga el flash sin abrir ning\'un stream de c\'amara ni
requerir el permiso \texttt{CAMERA} de manera obligatoria.

\subsection{B\'usqueda del ID de c\'amara con flash}

Cada dispositivo puede tener varias c\'amaras. Antes de usar la linterna hay
que identificar cu\'al tiene flash. Esto se resuelve en
\texttt{findTorchCameraId()}, llamada una \'unica vez desde \texttt{onCreate()}:

\begin{lstlisting}[caption={B\'usqueda del ID de c\'amara --- \texttt{MainActivity.java}}]
private String findTorchCameraId() {
    try {
        for (String id : cameraManager.getCameraIdList()) {
            Boolean hasFlash = cameraManager
                    .getCameraCharacteristics(id)
                    .get(CameraCharacteristics.FLASH_INFO_AVAILABLE);
            if (Boolean.TRUE.equals(hasFlash)) return id;
        }
    } catch (CameraAccessException e) {
        Log.e("MainActivity", "Impossible de trouver la camera torch", e);
    }
    return null;
}
\end{lstlisting}

\texttt{getCameraIdList()} devuelve un array con los identificadores de todas
las c\'amaras disponibles. Para cada una, se consulta la propiedad
\texttt{FLASH\_INFO\_AVAILABLE} de sus caracter\'isticas. Se retorna el primer
ID cuyo flash est\'e disponible, o \texttt{null} si ninguno lo tiene. Si el
resultado es \texttt{null}, \texttt{onShakeDetected()} no hace nada.

\subsection{Toggle de la linterna}

El encendido y apagado ocurre en \texttt{onShakeDetected()}, disparada por el
algoritmo de detecci\'on cuando se cumple el patr\'on de sacudida:

\begin{lstlisting}[caption={Toggle de la linterna --- \texttt{MainActivity.java}}]
private void onShakeDetected() {
    if (torchCameraId == null) return;
    try {
        cameraManager.setTorchMode(torchCameraId, !torchOn);
    } catch (CameraAccessException e) {
        Log.e("MainActivity", "Impossible de basculer la torche", e);
    }
}
\end{lstlisting}

Notar que \texttt{torchOn} \emph{no} se modifica aqu\'i directamente.
La fuente de verdad del estado real del flash es el \texttt{TorchCallback}
(secci\'on siguiente), lo que evita desincronizaciones si otra app o el
\texttt{TorchService} modifica el flash en paralelo.

\subsection{Sincronizaci\'on de estado con \texttt{TorchCallback}}

\texttt{CameraManager} provee un mecanismo de callback para recibir
notificaciones cuando el estado del flash cambia, sin importar el origen del
cambio:

\begin{lstlisting}[caption={Registro y uso del \texttt{TorchCallback} --- \texttt{MainActivity.java}}]
private final CameraManager.TorchCallback torchCallback =
        new CameraManager.TorchCallback() {
    @Override
    public void onTorchModeChanged(String cameraId, boolean enabled) {
        if (cameraId.equals(torchCameraId)) {
            torchOn = enabled;
            updateTorchUI();
        }
    }
};

// En onCreate():
cameraManager.registerTorchCallback(
        torchCallback,
        new Handler(Looper.getMainLooper()));

// En onDestroy():
cameraManager.unregisterTorchCallback(torchCallback);
\end{lstlisting}

El \texttt{Handler} con \texttt{Looper.getMainLooper()} garantiza que
\texttt{onTorchModeChanged()} se ejecute en el hilo principal, permitiendo
actualizar la UI directamente desde el callback sin necesidad de
\texttt{runOnUiThread()}.

\subsection{Actualizaci\'on del indicador visual}

Cuando el callback dispara \texttt{updateTorchUI()}, se modifica el color del
\texttt{View} circular (un \texttt{GradientDrawable}) y el texto de estado:

\begin{lstlisting}[caption={Actualizaci\'on del indicador de estado --- \texttt{MainActivity.java}}]
private void updateTorchUI() {
    GradientDrawable drawable =
            (GradientDrawable) viewTorchIndicator.getBackground();
    if (torchOn) {
        drawable.setColor(0xFFFFC107);      // ambar
        tvTorchStatus.setText("ALLUMÉE");
        tvTorchStatus.setTextColor(0xFFFFC107);
    } else {
        drawable.setColor(0xFF9E9E9E);      // gris
        tvTorchStatus.setText("ÉTEINTE");
        tvTorchStatus.setTextColor(0xFF9E9E9E);
    }
}
\end{lstlisting}

El color del drawable se modifica en tiempo de ejecuci\'on mediante
\texttt{setColor()}: el c\'irculo pasa de gris (\texttt{\#9E9E9E}) a \'ambar
(\texttt{\#FFC107}) al encender la linterna, y vuelve a gris al apagarla.
Esta es la \'unica funci\'on de la app que modifica elementos de UI, lo que
centraliza la l\'ogica de presentaci\'on en un \'unico lugar.

% ═════════════════════════════════════════════════════════════════════════════
\section{El Aceler\'ometro: \texttt{SensorManager}}
% ═════════════════════════════════════════════════════════════════════════════

\subsection{API utilizada}

El sistema de sensores de Android se accede a trav\'es de
\texttt{android.hardware.SensorManager}. Para recibir lecturas del
aceler\'ometro, \texttt{MainActivity} implementa la interfaz
\texttt{SensorEventListener}, que obliga a definir dos m\'etodos:
\texttt{onSensorChanged()} y \texttt{onAccuracyChanged()}.

El aceler\'ometro se obtiene e inicializa en \texttt{onCreate()}:

\begin{lstlisting}[caption={Inicializaci\'on del aceler\'ometro --- \texttt{MainActivity.java}}]
sensorManager = (SensorManager) getSystemService(Context.SENSOR_SERVICE);
if (sensorManager != null) {
    accelerometer = sensorManager.getDefaultSensor(Sensor.TYPE_ACCELEROMETER);
}
if (accelerometer == null) {
    Log.e("MainActivity", "Cet appareil n'a pas d'accelerometre");
}
\end{lstlisting}

\subsection{Ciclo de vida: registro y desregistro del listener}

Un aspecto cr\'itico del uso de sensores en Android es gestionar
correctamente el ciclo de vida. El listener se registra \'unicamente cuando
la app est\'a en primer plano y se desregistra al salir, lo que ahorra bater\'ia
y evita lecturas innecesarias:

\begin{lstlisting}[caption={Gesti\'on del ciclo de vida del sensor --- \texttt{MainActivity.java}}]
@Override
protected void onResume() {
    super.onResume();
    stopService(new Intent(this, TorchService.class));
    if (accelerometer != null) {
        sensorManager.registerListener(
                this,
                accelerometer,
                SensorManager.SENSOR_DELAY_NORMAL);
    }
}

@Override
protected void onPause() {
    super.onPause();
    if (accelerometer != null) {
        sensorManager.unregisterListener(this);
    }
    if (backgroundEnabled) {
        Intent intent = new Intent(this, TorchService.class);
        intent.putExtra(TorchService.EXTRA_THRESHOLD, shakeThreshold);
        intent.putExtra(TorchService.EXTRA_TORCH_STATE, torchOn);
        startForegroundService(intent);
    }
}
\end{lstlisting}

\texttt{SENSOR\_DELAY\_NORMAL} define una cadencia de muestreo de
aproximadamente 200\,ms entre lecturas. En \texttt{onResume()} tambi\'en se
detiene el \texttt{TorchService} si estuviese corriendo, para garantizar que
en todo momento \emph{solo un} componente escucha el sensor.

% ═════════════════════════════════════════════════════════════════════════════
\section{Detecci\'on de Sacudida}
% ═════════════════════════════════════════════════════════════════════════════

\subsection{C\'alculo de la magnitud}

El aceler\'ometro entrega tres valores: aceleraci\'on en los ejes $X$, $Y$ y
$Z$ expresada en m/s$^2$. Para detectar una sacudida independientemente de la
orientaci\'on del tel\'efono, se calcula la \textbf{magnitud del vector de
aceleraci\'on} normalizada en unidades de gravedad (g-force):

\[
  g_{\text{force}} =
    \sqrt{
      \left(\frac{a_x}{g}\right)^2 +
      \left(\frac{a_y}{g}\right)^2 +
      \left(\frac{a_z}{g}\right)^2
    }
\]

donde $g = 9{,}80665\,\text{m/s}^2$ (\texttt{SensorManager.GRAVITY\_EARTH}).
Con el tel\'efono en reposo, el vector apunta en la direcci\'on de la gravedad
con magnitud $\approx 1\,g$. Una sacudida produce picos de $2{,}5\,g$--$5\,g$.

\subsection{Algoritmo de detecci\'on}

La funci\'on \texttt{onSensorChanged()} se ejecuta con cada nueva lectura del
sensor y concentra toda la l\'ogica de detecci\'on:

\begin{lstlisting}[caption={Detecci\'on de sacudida en tiempo real --- \texttt{MainActivity.java}}]
@Override
public void onSensorChanged(SensorEvent event) {
    float gX = event.values[0] / SensorManager.GRAVITY_EARTH;
    float gY = event.values[1] / SensorManager.GRAVITY_EARTH;
    float gZ = event.values[2] / SensorManager.GRAVITY_EARTH;
    float gForce = (float) Math.sqrt(gX * gX + gY * gY + gZ * gZ);

    graphAccelerometer.addValue(gForce);   // actualiza el grafico

    if (gForce < shakeThreshold) return;   // filtro 1: umbral

    long now = System.currentTimeMillis();
    if (now - lastShakeTime < SHAKE_SLOP_TIME_MS) return; // filtro 2

    if (now - lastShakeTime > SHAKE_COUNT_RESET_TIME_MS) {
        shakeCount = 0;                    // filtro 3: ventana temporal
    }

    lastShakeTime = now;
    shakeCount++;

    if (shakeCount >= MIN_SHAKE_COUNT) {
        shakeCount = 0;
        onShakeDetected();
    }
}
\end{lstlisting}

El algoritmo funciona en cuatro capas de filtrado superpuestas:

\begin{enumerate}
  \item \textbf{Filtro de umbral} (\texttt{shakeThreshold}): si
        \texttt{gForce} no supera el umbral, la lectura se descarta. El umbral
        es ajustado en tiempo real por el slider de sensibilidad (de
        $4{,}0\,g$ poco sensible a $1{,}5\,g$ muy sensible).

  \item \textbf{Tiempo m\'inimo entre sacudidas}
        (\texttt{SHAKE\_SLOP\_TIME\_MS = 500\,ms}): evita que un \'unico
        movimiento brusco cuente como m\'ultiples sacudidas mientras la
        lectura del sensor sigue superando el umbral.

  \item \textbf{Ventana temporal de reinicio}
        (\texttt{SHAKE\_COUNT\_RESET\_TIME\_MS = 3000\,ms}): si transcurren
        m\'as de 3 segundos entre dos sacudidas v\'alidas, el contador se pone
        a cero. Esto obliga a que las sacudidas sean consecutivas y no
        acumuladas a lo largo del tiempo.

  \item \textbf{Contador m\'inimo} (\texttt{MIN\_SHAKE\_COUNT = 2}): se
        necesitan al menos 2 sacudidas v\'alidas y consecutivas para disparar
        \texttt{onShakeDetected()}. Esto previene activaciones accidentales por
        un \'unico golpe.
\end{enumerate}

\subsection{Constantes y umbral configurable}

Los par\'ametros de tiempo son constantes de clase; \texttt{shakeThreshold}
es un campo mutable para que el slider pueda modificarlo en tiempo real:

\begin{lstlisting}[caption={Constantes del algoritmo de detecci\'on --- \texttt{MainActivity.java}}]
private static final int SHAKE_SLOP_TIME_MS        = 500;
private static final int SHAKE_COUNT_RESET_TIME_MS = 3000;
private static final int MIN_SHAKE_COUNT           = 2;

private float shakeThreshold = 2.75f;
\end{lstlisting}

La f\'ormula de mapeo del slider al umbral:

\begin{lstlisting}[caption={Ajuste din\'amico del umbral desde el \texttt{SeekBar} --- \texttt{MainActivity.java}}]
seekBarSensitivity.setOnSeekBarChangeListener(
        new SeekBar.OnSeekBarChangeListener() {
    @Override
    public void onProgressChanged(SeekBar seekBar,
                                  int progress,
                                  boolean fromUser) {
        shakeThreshold = 4.0f - (progress * 0.25f);
        tvSensitivityValue.setText(progress + " / 10");
        graphAccelerometer.setThreshold(shakeThreshold);
    }
    @Override public void onStartTrackingTouch(SeekBar seekBar) {}
    @Override public void onStopTrackingTouch(SeekBar seekBar) {}
});
\end{lstlisting}

La f\'ormula \texttt{4.0 - progress * 0.25} mapea el rango $[0,\,10]$ del
slider al rango de umbral $[4{,}0\,g,\;1{,}5\,g]$: a mayor nivel de
sensibilidad (slider a la derecha), menor es el umbral y m\'as f\'acil resulta
disparar la detecci\'on. El valor por defecto es \texttt{progress = 5}
($\rightarrow 2{,}75\,g$).

% ═════════════════════════════════════════════════════════════════════════════
\section{Gr\'afico en Tiempo Real: \texttt{AccelerometerGraphView}}
% ═════════════════════════════════════════════════════════════════════════════

\subsection{Vista personalizada con Canvas}

Para visualizar los valores del aceler\'ometro en tiempo real se implement\'o
\texttt{AccelerometerGraphView}, una clase que extiende
\texttt{android.view.View} y dibuja directamente sobre un \texttt{Canvas}.

El gr\'afico muestra:
\begin{itemize}
  \item \textbf{L\'inea azul}: magnitud \texttt{gForce} en tiempo real.
  \item \textbf{L\'inea roja punteada}: umbral actual, se desplaza
        verticalmente al mover el slider.
  \item \textbf{Grilla gris}: l\'ineas horizontales cada 1\,g.
\end{itemize}

\subsection{Buffer circular}

La vista mantiene un buffer circular de 150 muestras. Al llegar una nueva
lectura, sobreescribe la posici\'on m\'as antigua y el puntero \texttt{head}
avanza en m\'odulo:

\begin{lstlisting}[caption={Buffer circular y actualizaci\'on --- \texttt{AccelerometerGraphView.java}}]
private static final int   BUFFER_SIZE = 150;
private static final float Y_MAX       = 5.0f;

private final float[] values = new float[BUFFER_SIZE];
private int   head      = 0;
private int   count     = 0;
private float threshold = 2.75f;

public void addValue(float gForce) {
    values[head] = gForce;
    head = (head + 1) % BUFFER_SIZE;
    if (count < BUFFER_SIZE) count++;
    postInvalidate();
}

public void setThreshold(float threshold) {
    this.threshold = threshold;
    postInvalidate();
}
\end{lstlisting}

\texttt{postInvalidate()} es \emph{thread-safe}: puede llamarse desde cualquier
hilo (incluyendo el hilo del sensor) y programa un redibujo en el hilo de UI.
El eje Y del gr\'afico cubre de $0$ a $5\,g$ (\texttt{Y\_MAX}), rango que
contiene en todo momento la l\'inea de umbral (entre $1{,}5\,g$ y $4{,}0\,g$).

\subsection{Inicializaci\'on de los pinceles}

Los objetos \texttt{Paint} se crean \emph{\'unicamente en el constructor}, nunca
dentro de \texttt{onDraw()}. Esta es una regla fundamental en Android: instanciar
objetos en \texttt{onDraw()} genera basura en cada frame y degrada
significativamente el rendimiento.

\begin{lstlisting}[caption={Inicializaci\'on de \texttt{Paint} en el constructor --- \texttt{AccelerometerGraphView.java}}]
public AccelerometerGraphView(Context context, AttributeSet attrs) {
    super(context, attrs);

    linePaint = new Paint(Paint.ANTI_ALIAS_FLAG);
    linePaint.setColor(0xFF42A5F5);      // azul Material Blue 400
    linePaint.setStrokeWidth(3f);
    linePaint.setStyle(Paint.Style.STROKE);

    thresholdPaint = new Paint(Paint.ANTI_ALIAS_FLAG);
    thresholdPaint.setColor(0xFFEF5350); // rojo Material Red 400
    thresholdPaint.setStrokeWidth(2f);
    thresholdPaint.setStyle(Paint.Style.STROKE);
    thresholdPaint.setPathEffect(
            new DashPathEffect(new float[]{16f, 8f}, 0));

    gridPaint = new Paint();
    gridPaint.setColor(0x33FFFFFF);      // blanco al 20% de opacidad
    gridPaint.setStrokeWidth(1f);
}
\end{lstlisting}

El \texttt{DashPathEffect} de la l\'inea de umbral recibe un array de dos
valores: $[16, 8]$ significa segmento de 16\,px encendido, 8\,px apagado,
generando el efecto de l\'inea punteada.

\subsection{M\'etodo \texttt{onDraw}}

\texttt{onDraw()} se ejecuta cada vez que la vista necesita redibujarse. Dibuja
tres capas en orden sobre el \texttt{Canvas}:

\begin{lstlisting}[caption={Dibujo del gr\'afico con Canvas --- \texttt{AccelerometerGraphView.java}}]
@Override
protected void onDraw(Canvas canvas) {
    int w = getWidth();
    int h = getHeight();

    canvas.drawColor(0xFF1E1E1E);           // fondo oscuro VSCode

    for (int g = 1; g <= 4; g++) {         // grilla cada 1g
        float y = h * (1f - g / Y_MAX);
        canvas.drawLine(0, y, w, y, gridPaint);
    }

    float ty = h * (1f - threshold / Y_MAX); // linea de umbral
    canvas.drawLine(0, ty, w, ty, thresholdPaint);

    if (count < 2) return;

    path.reset();
    float xStep = (float) w / (BUFFER_SIZE - 1);
    for (int i = 0; i < count; i++) {
        int   idx = (head - count + i + BUFFER_SIZE) % BUFFER_SIZE;
        float x   = i * xStep;
        float y   = h * (1f - Math.min(values[idx], Y_MAX) / Y_MAX);
        if (i == 0) path.moveTo(x, y);
        else        path.lineTo(x, y);
    }
    canvas.drawPath(path, linePaint);
}
\end{lstlisting}

La coordenada vertical en p\'ixeles para un valor $v$ se calcula como:

\[
  y_{\text{px}} = h \cdot \left(1 - \frac{v}{Y_{\max}}\right)
\]

El factor $(1 - \cdot)$ invierte el eje porque en el sistema de coordenadas
del Canvas de Android el eje $Y$ crece hacia abajo (el origen est\'a en la
esquina superior izquierda). El \'indice \texttt{idx} recorre el buffer
circular en orden cronol\'ogico: del dato m\'as antiguo (izquierda) al m\'as
reciente (derecha), produciendo el efecto de scroll hacia la izquierda a
medida que llegan nuevas muestras.

% ═════════════════════════════════════════════════════════════════════════════
\section{Ejecuci\'on en Segundo Plano: \texttt{TorchService}}
% ═════════════════════════════════════════════════════════════════════════════

\subsection{Motivaci\'on}

El comportamiento predeterminado de Android es suspender el aceler\'ometro
cuando la Activity entra en pausa (al salir de la app). Para mantener la
funcionalidad activa en segundo plano, Android exige usar un
\textbf{Foreground Service}: un servicio que muestra una notificaci\'on
persistente, informando al usuario que la aplicaci\'on sigue activa. Esto
impide que el sistema operativo mate el proceso por falta de recursos.

\subsection{Arquitectura del servicio}

\texttt{TorchService} extiende \texttt{Service} e implementa
\texttt{SensorEventListener}, replicando exactamente la misma l\'ogica de
detecci\'on de sacudidas de \texttt{MainActivity}:

\begin{lstlisting}[caption={Arranque del servicio --- \texttt{TorchService.java}}]
@Override
public int onStartCommand(Intent intent, int flags, int startId) {
    if (intent != null) {
        shakeThreshold = intent.getFloatExtra(EXTRA_THRESHOLD, 2.75f);
        torchOn        = intent.getBooleanExtra(EXTRA_TORCH_STATE, false);
    }
    createNotificationChannel();
    startForeground(NOTIF_ID, buildNotification());
    if (accelerometer != null) {
        sensorManager.registerListener(
                this,
                accelerometer,
                SensorManager.SENSOR_DELAY_NORMAL);
    }
    return START_STICKY;
}
\end{lstlisting}

El servicio recibe el umbral de sensibilidad y el estado actual del flash como
\texttt{Intent extras}, para continuar desde donde lo dej\'o la Activity.
\texttt{startForeground()} debe invocarse antes de que transcurran 5 segundos
desde el inicio del servicio. \texttt{START\_STICKY} indica al sistema operativo
que, si necesita terminar el proceso, debe reiniciarlo autom\'aticamente.

La notificaci\'on persistente se construye con \texttt{NotificationCompat}:

\begin{lstlisting}[caption={Construcci\'on de la notificaci\'on --- \texttt{TorchService.java}}]
private Notification buildNotification() {
    Intent intent = new Intent(this, MainActivity.class);
    PendingIntent pi = PendingIntent.getActivity(
            this, 0, intent, PendingIntent.FLAG_IMMUTABLE);
    return new NotificationCompat.Builder(this, CHANNEL_ID)
            .setContentTitle("Helios actif")
            .setContentText("Secouez pour allumer/eteindre la torche")
            .setSmallIcon(R.mipmap.ic_launcher)
            .setContentIntent(pi)
            .build();
}
\end{lstlisting}

El \texttt{PendingIntent} con \texttt{FLAG\_IMMUTABLE} permite al usuario
tocar la notificaci\'on para volver a la Activity.

\subsection{Traspaso de control entre Activity y Service}

El dise\~no garantiza que en todo momento \emph{solo uno} de los dos
componentes escucha el aceler\'ometro:

\begin{center}
\begin{tabular}{@{} l l l @{}}
  \toprule
  \textbf{Estado de la app} & \textbf{Toggle fondo} & \textbf{Qui\'en escucha} \\
  \midrule
  En primer plano & OFF & MainActivity \\
  En primer plano & ON  & MainActivity \\
  En segundo plano & OFF & Nadie (sensor inactivo) \\
  En segundo plano & ON  & TorchService \\
  \bottomrule
\end{tabular}
\end{center}

Al volver a la app, \texttt{onResume()} detiene el servicio \emph{antes} de
registrar el listener en la Activity:

\begin{lstlisting}[caption={Traspaso de vuelta a la Activity --- \texttt{MainActivity.java}}]
@Override
protected void onResume() {
    super.onResume();
    stopService(new Intent(this, TorchService.class));
    if (accelerometer != null) {
        sensorManager.registerListener(
                this,
                accelerometer,
                SensorManager.SENSOR_DELAY_NORMAL);
    }
}
\end{lstlisting}

Al salir de la app, \texttt{onPause()} desregistra el listener de la Activity
y, si el toggle est\'a activo, lanza el servicio:

\begin{lstlisting}[caption={Traspaso hacia el servicio al pausar --- \texttt{MainActivity.java}}]
@Override
protected void onPause() {
    super.onPause();
    if (accelerometer != null) {
        sensorManager.unregisterListener(this);
    }
    if (backgroundEnabled) {
        Intent intent = new Intent(this, TorchService.class);
        intent.putExtra(TorchService.EXTRA_THRESHOLD, shakeThreshold);
        intent.putExtra(TorchService.EXTRA_TORCH_STATE, torchOn);
        startForegroundService(intent);
    }
}
\end{lstlisting}

\subsection{Toggle en la interfaz}

El usuario controla el modo segundo plano con un \texttt{SwitchMaterial}
ubicado en la esquina superior derecha de la pantalla. Su estado se almacena
en el campo \texttt{backgroundEnabled}:

\begin{lstlisting}[caption={Toggle de modo segundo plano --- \texttt{MainActivity.java}}]
SwitchMaterial switchBackground = findViewById(R.id.switchBackground);
switchBackground.setOnCheckedChangeListener(
        (buttonView, isChecked) -> backgroundEnabled = isChecked);
\end{lstlisting}

\subsection{Declaraci\'on en el Manifest}

Para que Android reconozca el servicio y le permita ejecutarse en primer
plano, se lo declara en \texttt{AndroidManifest.xml}:

\begin{lstlisting}[style=vsxml, caption={Permisos y declaraci\'on del servicio --- \texttt{AndroidManifest.xml}}]
<uses-permission android:name="android.permission.CAMERA" />
<uses-permission android:name="android.permission.FOREGROUND_SERVICE" />
<uses-permission android:name=
    "android.permission.FOREGROUND_SERVICE_SPECIAL_USE" />
<uses-permission android:name="android.permission.POST_NOTIFICATIONS" />

<service
    android:name=".TorchService"
    android:foregroundServiceType="specialUse"
    android:exported="false" />
\end{lstlisting}

El tipo \texttt{specialUse} es el m\'as apropiado para este caso: la
combinaci\'on de aceler\'ometro y control de flash no encaja en ninguna de las
categor\'ias espec\'ificas de Android~14 (c\'amara, ubicaci\'on, salud, etc.).
\texttt{POST\_NOTIFICATIONS} es requerido para mostrar la notificaci\'on
persistente en Android~13 y superior. El permiso se solicita en tiempo de
ejecuci\'on desde \texttt{onCreate()} de la Activity.

% ═════════════════════════════════════════════════════════════════════════════
\section{Conclusiones}
% ═════════════════════════════════════════════════════════════════════════════

Android Helios integra los dos m\'odulos de hardware requeridos---aceler\'ometro
y c\'amara/flash---de forma funcional y robusta. Los puntos clave del dise\~no son:

\begin{itemize}
  \item \textbf{Detecci\'on de sacudidas en cuatro capas}: umbral de g-force
        $+$ tiempo m\'inimo entre sacudidas $+$ ventana de reinicio $+$
        contador m\'inimo. Cada capa reduce los falsos positivos
        independientemente.

  \item \textbf{TorchCallback como fuente \'unica de verdad}: el estado de la
        linterna siempre refleja el estado real del hardware, sin importar
        qu\'e componente lo cambi\'o (Activity, Service, u otra app).

  \item \textbf{Gr\'afico como herramienta de calibraci\'on}: el usuario puede
        ver en tiempo real la magnitud de sus movimientos y ajustar el umbral
        hasta que la l\'inea roja quede por debajo de sus sacudidas y por
        encima del ruido de fondo.

  \item \textbf{Ciclo de vida estricto}: el sensor nunca corre en segundo
        plano sin consentimiento expl\'icito, y cuando lo hace, el sistema
        operativo es notificado correctamente a trav\'es del Foreground
        Service con su notificaci\'on persistente.
\end{itemize}

\end{document}
